\chapter{Review of Algebraic Prerequisites}
\label{app:algebra}

This appendix collects, without proof, the algebraic facts assumed throughout
the text: properties of exponents and radicals, factoring techniques, the
quadratic formula, and the algebra of rational expressions. Students who find
early exercises unexpectedly difficult are encouraged to consult it first.

\section{Exponents and Radicals}

For real numbers $a, b$ and integers $m, n$ (with $a, b \neq 0$ where
division or a nonpositive exponent is involved), the exponent rules are
\[
a^m a^n = a^{m+n}, \qquad \frac{a^m}{a^n} = a^{m-n}, \qquad (a^m)^n =
a^{mn}, \qquad (ab)^n = a^n b^n, \qquad a^0 = 1, \qquad a^{-n} =
\frac{1}{a^n}.
\]
A radical is related to a fractional exponent by $\sqrt[n]{a} =
a^{1/n}$, and more generally $a^{m/n} = \sqrt[n]{a^m} =
\left(\sqrt[n]{a}\right)^m$, provided $\sqrt[n]{a}$ is defined. The
following radical identities hold whenever both sides are defined:
\[
\sqrt[n]{ab} = \sqrt[n]{a}\,\sqrt[n]{b}, \qquad
\sqrt[n]{\frac{a}{b}} = \frac{\sqrt[n]{a}}{\sqrt[n]{b}}, \qquad
\sqrt[n]{a^n} = |a| \text{ when } n \text{ is even}, \qquad
\sqrt[n]{a^n} = a \text{ when } n \text{ is odd}.
\]
A radical in a denominator is customarily removed by rationalizing:
multiplying numerator and denominator by $\sqrt{a}$ clears a single
square root, as in $\dfrac{1}{\sqrt{a}} = \dfrac{\sqrt{a}}{a}$, while a
denominator of the form $a + \sqrt{b}$ is cleared by multiplying by its
conjugate $a - \sqrt{b}$, using the difference-of-squares identity of the
next section.

\section{Factoring}

The following factoring patterns recur throughout the text.
\begin{align*}
\text{Common factor:} \qquad & ax + ay = a(x+y). \\
\text{Difference of squares:} \qquad & a^2 - b^2 = (a-b)(a+b). \\
\text{Perfect square trinomial:} \qquad & a^2 \pm 2ab + b^2 = (a \pm
b)^2. \\
\text{Difference of cubes:} \qquad & a^3 - b^3 = (a-b)(a^2+ab+b^2). \\
\text{Sum of cubes:} \qquad & a^3 + b^3 = (a+b)(a^2-ab+b^2). \\
\text{Trinomial:} \qquad & x^2 + (p+q)x + pq = (x+p)(x+q).
\end{align*}
For a general quadratic $ax^2+bx+c$ with $a \neq 1$, factoring by
grouping is often effective: rewrite the middle term as a sum of two
terms whose coefficients multiply to $ac$ and add to $b$, then factor the
resulting four-term expression in pairs.

The \emph{quadratic formula} solves $ax^2+bx+c=0$ for $a \neq 0$:
\[
x = \frac{-b \pm \sqrt{b^2-4ac}}{2a}.
\]
The quantity $b^2-4ac$, called the \emph{discriminant}, determines the
number and type of solutions: two distinct real solutions when it is
positive, one repeated real solution when it is zero, and two complex
conjugate solutions when it is negative, as discussed in Chapter~16.

\section{Solving Equations}

An equation is solved by applying the same operation to both sides until
the unknown is isolated, using operations that preserve the solution set:
adding or subtracting the same quantity from both sides, multiplying or
dividing both sides by the same nonzero quantity, and factoring a side
that equals zero. Two operations deserve particular caution, both
discussed in the context of trigonometric equations in Chapter~11:
squaring both sides of an equation can introduce extraneous solutions
not present in the original equation, and multiplying both sides by an
expression that could equal zero, or that involves the unknown, can
likewise change the solution set. Every candidate solution obtained by
such an operation should be checked in the original equation before being
accepted. A rational equation, one containing the unknown in a
denominator, is typically solved by multiplying through by the least
common denominator to clear the fractions, solving the resulting
polynomial equation, and then discarding any solution that makes an
original denominator equal to zero.
